\documentclass[../main.tex]{subfiles}

\begin{document}

\chapter{实数集：集合、序列与函数}\label{chapter:ch01-实数集}

假定读者已经熟悉实数、实数集、实数序列、一元实变量实值函数的性质, 这些通常在本科生的分析课程中讨论. 具有这些背景知识将使得读者能够透彻理解本章, 本章用于快速而全面地建立以后要使用和参考的结果. 假定实数集, 记为 $\mathbf{R}$, 满足三种类型的公理. 我们叙述这些公理并从中导出自然数、有理数以及不可数集的性质. 有了这些作为背景, 我们建立实数的开集与闭集的性质, 收敛、单调、实数的 Cauchy 序列, 一元实变量连续实值函数.

\subfile{Chapter-01/section-01.tex}

\subfile{Chapter-01/section-02.tex}

\subfile{Chapter-01/section-03.tex}

\subfile{Chapter-01/section-04.tex}

\subfile{Chapter-01/section-05.tex}

\subfile{Chapter-01/section-06.tex}

\end{document}